\documentclass{ctexart}
\usepackage{amsmath, amssymb}
\usepackage{geometry}
\geometry{a4paper, margin=2cm}

\title{\textbf{第8章 机械波} --- 知识点与公式}
\author{}
\date{}

\begin{document}
\maketitle

\section*{8.1 机械波的产生和传播}

产生条件：波源 + 弹性介质。

横波：质点振动方向 $\perp$ 波传播方向。
纵波：质点振动方向 $\parallel$ 波传播方向。

波面（波阵面）：振动相位相同的点构成的面。
波线：沿波传播方向的有向线。
波前：某一时刻波传播到的最前面的波面。

波长 $\lambda$：同一波线上相邻相位差 $2\pi$ 的两点间距离。
周期 $T$：波前进一个波长所需时间。
频率 $\nu = 1/T$。
波速 $u = \dfrac{\lambda}{T} = \nu\lambda$。

波速公式：
弦线横波：$u = \sqrt{\dfrac{T}{\mu}}$（$T$ 张力，$\mu$ 线密度）
固体棒纵波：$u = \sqrt{\dfrac{Y}{\rho}}$（$Y$ 杨氏模量）

\section*{8.2 平面简谐波}

波函数（沿 $x$ 正方向传播）：
\[
y = A\cos\left[\omega\left(t - \frac{x}{u}\right) + \varphi_0\right]
= A\cos\left[2\pi\left(\frac{t}{T} - \frac{x}{\lambda}\right) + \varphi_0\right]
= A\cos(\omega t - kx + \varphi_0)
\]

沿 $x$ 负方向传播：$y = A\cos\left[\omega\left(t + \dfrac{x}{u}\right) + \varphi_0\right]$

波数 $k = \dfrac{2\pi}{\lambda}$

\subsection*{波函数的物理意义}
\begin{itemize}
  \item $x$ 给定：$y = y(t)$ 为该点的振动方程
  \item $t$ 给定：$y = y(x)$ 为 $t$ 时刻的波形图
  \item $x, t$ 均变化：波形传播，体现时空周期性
\end{itemize}

\section*{8.3 波的能量}

能量密度：$w = \dfrac{\mathrm{d}W}{\mathrm{d}V} = \rho A^{2}\omega^{2}\sin^{2}\left[\omega\left(t - \dfrac{x}{u}\right) + \varphi_0\right]$

平均能量密度：$\bar{w} = \dfrac12 \rho A^{2}\omega^{2}$

平均能流密度（波的强度）：$I = \bar{w}u = \dfrac12 \rho A^{2}\omega^{2}u$

\section*{8.5 波的干涉}

相干条件：频率相同、振动方向相同、相位差恒定。

两列波在 $P$ 点的相位差：
\[
\Delta\varphi = \varphi_2 - \varphi_1 - 2\pi\frac{r_2 - r_1}{\lambda}
\]

干涉加强条件：$\Delta\varphi = \pm 2k\pi$，$A = A_1 + A_2$
干涉减弱条件：$\Delta\varphi = \pm (2k+1)\pi$，$A = |A_1 - A_2|$

\section*{8.6 驻波}

驻波方程：
\[
y = 2A\cos\frac{2\pi x}{\lambda}\cos\omega t
\]

波腹位置：$x = \pm k\dfrac{\lambda}{2}$（$k = 0,1,2,\dots$）
波节位置：$x = \pm (2k+1)\dfrac{\lambda}{4}$（$k = 0,1,2,\dots$）

相邻波节（或波腹）间距：$\dfrac{\lambda}{2}$

弦线驻波条件：$L = n\dfrac{\lambda}{2}$，$\nu = n\dfrac{u}{2L}$（$n = 1,2,3,\dots$）

半波损失：波从波疏介质入射到波密介质界面反射时，反射波有 $\pi$ 相位突变。

\section*{8.7 多普勒效应}

观察者接收频率：
\[
\nu' = \frac{u + v_0}{u - v_s}\nu_0
\]

$v_0$：观察者速度（靠近波源为正），$v_s$：波源速度（靠近观察者为正）。

电磁波多普勒效应：
\[
\nu' = \sqrt{\frac{1 \pm v/c}{1 \mp v/c}}\,\nu_s
\]

冲击波马赫角：$\sin\theta = \dfrac{u}{v_s}$（$v_s > u$）

\end{document}
